\section{Introduction}
\label{sec:intro}

Modern production systems increasingly rely on machine learning for operational intelligence. Anomaly detection identifies performance regressions before users are impacted, predictive models forecast resource bottlenecks, and automated root cause analysis reduces mean time to resolution.
These ML-driven observability tools share a common dependency: they require large volumes of representative execution traces for training.
At scale, kernel-level traces provide the most comprehensive view of system behavior, capturing scheduler decisions, memory allocations, I/O operations, and inter-process communication with microsecond-level precision.

\noindent\textbf{The Data Scarcity Problem.}
Collecting production kernel traces is expensive and constrained by multiple factors.
First, kernel tracing imposes non-trivial runtime overhead (typically 2--5\% CPU utilization with tools such as LTTng~\cite{Desnoyers2006The}), making continuous collection impractical for latency-sensitive services.
Second, kernel traces often contain sensitive information---including process identifiers, file paths, and network endpoints---that may violate data retention policies or require costly sanitization before use in ML pipelines.
Third, and most critically, real traces capture only behaviors that occurred during limited collection windows: rare failure modes, edge-case workloads, and uncommon scheduling patterns remain underrepresented.
This long-tail problem is particularly acute for ML-based diagnostics, where the events of interest are precisely those least likely to appear during normal operation.

As a result, practitioners face a fundamental tension: ML models require diverse and representative training data, yet production constraints limit what can be safely and cost-effectively collected.
In practice, teams often resort to opportunistic trace collection during maintenance windows or from staging environments, yielding datasets that poorly reflect production diversity.

\noindent\textbf{Limitations of Existing Approaches.}
Prior work on trace synthesis has largely focused on application-level logs~\cite{logsynth} or network traffic~\cite{netgan}, which lack the fine-grained system state captured in kernel traces.
Statistical models such as Markov chains can generate locally valid event sequences but fail to capture long-range temporal dependencies and multi-attribute correlations inherent in kernel execution.
Rule-based generators require substantial domain expertise and do not generalize easily across workloads.
Overall, existing approaches do not adequately address the challenge of generating synthetic kernel traces that simultaneously preserve temporal structure, multi-channel correlations, and system-level invariants while remaining practical for ML pipeline integration.

\noindent\textbf{Our Approach.}
We present \textsc{SynthTrace}, a framework for generating synthetic kernel execution traces using diffusion models.
\textsc{SynthTrace} represents kernel traces as multi-channel sequences, jointly modeling event types, inter-event timing, CPU affinity, thread identifiers, process names, and return values.
A Transformer-based denoising diffusion model learns to generate coherent execution traces by iteratively refining noise into structured sequences over configurable temporal context lengths.
To mitigate violations of system-level invariants—such as invalid event transitions, implausible timing relationships, or inconsistent attribute values—we introduce a constraint learning module that mines invariants from real traces and a post-hoc repair mechanism that detects and corrects violations without retraining the generative model.

\noindent\textbf{Contributions.}
This paper makes the following contributions:

\begin{itemize}[leftmargin=*,nosep]
  \item \textbf{C1: Problem characterization.} We characterize the data scarcity challenge for ML-based system diagnostics and identify key requirements for synthetic kernel trace generation, including multi-channel fidelity, long-range temporal coherence, and constraint awareness.

  \item \textbf{C2: \textsc{SynthTrace} framework.} We propose a diffusion-based trace generator that jointly models six kernel trace channels and supports context lengths up to 4096 events. The framework incorporates learned execution invariants and a constraint-guided repair mechanism to improve system-level validity.

\item \textbf{C3: Empirical evaluation of utility and safety.}
We evaluate \textsc{SynthTrace} on diverse benchmark workloads using downstream next-event prediction as a proxy for trace utility. Our results show that synthetic traces can effectively augment limited real data for structured, compute-heavy workloads, while exhibiting workload-dependent limitations for I/O-intensive systems. Constraint-guided repair consistently mitigates violation-induced degradation with low risk.


  \item \textbf{C4: Practical guidance.} We provide actionable insights on context length selection, feature complexity trade-offs, and when synthetic kernel traces can be safely integrated into ML pipelines.
\end{itemize}

\noindent\textbf{Paper Organization.}
The remainder of this paper is structured as follows.
Section~\ref{sec:background} provides background on diffusion models and related work.
Section~\ref{sec:approach} describes the \textsc{SynthTrace} framework and constraint-guided repair mechanism.
Section~\ref{sec:evaluation} details the experimental setup and evaluation methodology.
Section~\ref{sec:results} presents results addressing our four research questions.
Section~\ref{sec:discussion} synthesizes findings and discusses practical implications.
Finally, Section~\ref{sec:limitations-and-threats-to-validity} outlines limitations and threats to validity.


% Machine learning models for system behavior analysis require large volumes of training data, yet collecting comprehensive kernel execution traces from production systems is expensive, time-consuming, and often infeasible due to privacy constraints or instrumentation overhead. Synthetic data generation offers a promising alternative: train generative models on existing traces to produce realistic synthetic sequences that can augment or replace real data.

% Recent advances in diffusion probabilistic models have demonstrated success in generating high-fidelity images, text, and audio. However, kernel execution traces present unique challenges. Unlike natural data, kernel traces are multi-modal sequences with strict structural constraints: event transitions must follow valid system call orderings, CPU affinity must respect scheduling rules, and timing patterns must reflect realistic system dynamics. Violating these constraints produces invalid traces unsuitable for applications requiring hard guarantees of system validity, such as trace replay or deterministic debugging.

% This paper presents a diffusion-based approach for generating synthetic kernel execution traces. We develop a multi-channel diffusion model that jointly generates six correlated trace attributes (event type, inter-event time, CPU core, thread ID, process command, return value) in a unified latent space. To ensure structural validity, we introduce a constraint-guided repair mechanism that learns valid patterns from real traces and applies post-hoc corrections. We evaluate our approach on nine diverse Phoronix benchmarks using a next-event prediction task as a proxy for synthetic data quality.

% Our evaluation reveals that synthetic data quality is fundamentally workload-dependent. For deterministic, structured workloads (e.g., numerical computation), synthetic augmentation achieves near-parity with real data, with only 2.6\% F1-Macro degradation at context length L=4096. In contrast, non-deterministic I/O-heavy workloads show significant degradation (25--36\%), indicating fundamental challenges for current generative approaches. Context length emerges as the dominant factor for quality, with L=4096 doubling average performance compared to L=256. Ablation studies reveal that simpler 2-channel models achieve comparable quality to full 6-channel models with significantly lower computational cost.

% \noindent\textbf{Contributions.} This work makes the following contributions:

% \begin{itemize}[leftmargin=*,nosep]
%   \item A multi-channel diffusion model for kernel trace generation that jointly models event semantics, timing patterns, and system resource allocation.

%   \item A constraint-guided repair mechanism that ensures 100\% structural validity while preserving empirical distributions of real traces.

%   \item A comprehensive evaluation demonstrating that synthetic data achieves near-parity with real data for deterministic workloads, with actionable guidance on when synthetic augmentation is viable.

%   \item Systematic ablation studies identifying optimal context length (L=4096) and cost-effective model configurations (2-channel models suffice for most workloads).
% \end{itemize}

% \noindent\textbf{Paper Organization.} The remainder of this paper is structured as follows: Section~\ref{sec:background} provides background on diffusion models and related work on trace generation. Section~\ref{sec:approach} describes our multi-channel diffusion model, constraint-guided repair mechanism, and evaluation metrics. Section~\ref{sec:experiments} details the experimental setup including downstream task configuration and training procedures. Section~\ref{sec:results} presents results for our four research questions. Section~\ref{sec:discussion} synthesizes findings and provides actionable insights for practitioners. Finally, Section~\ref{sec:limitations-and-threats-to-validity} discusses threats to validity and limitations